\section{Permanent-specific containment and the remaining rank comparison}

\subsection*{Finding F-009: a permanent-size gate valid for every partition}
The following uses a known upper bound.  Grenet,
\emph{An upper bound for the permanent versus determinant problem},
25 May 2012 version, unnumbered Lemma on p.~2 and unnumbered
Theorem on p.~3, constructs an affine determinantal representation
of size \(2^m-1\).  We record the homogeneous identity and its
coordinate-ring consequence explicitly.

Let \(m\geq2\), put \(s=2^m-1\), and first let \(L\) be an
auxiliary indeterminate.  There is an \(s\times s\) matrix of homogeneous
linear forms over \(\mathbb Z\) such that
\begin{equation}\label{eq:F009}
                    \det M=L^{s-m}\operatorname{perm}_m(X).
\end{equation}
Consequently, for every \(n\geq2^m-1\), the entire padded-permanent
orbit closure lies in the determinant orbit closure:
\[
                         Z_{n,m}\subseteq\Omega_n.
\]
In particular \(q^m_\lambda(d[n])\leq K_n(\lambda)\) for every
\(d\) and every partition \(\lambda\vdash nd\), without a row bound.

\begin{proof}[Proof of the homogeneous form and consequence]
Use one vertex for each subset of \([m]\), identifying the empty
and full subsets as a central vertex.  There are \(s\) vertices.
For \(S\subsetneq[m]\) and \(j\notin S\), put an edge
\(S\to S\cup\{j\}\), of weight \(x_{|S|+1,j}\).
Multiply the weights of all edges leaving the central vertex
by \((-1)^{m-1}\).  Put a loop of weight \(L\) at every
noncentral vertex and no loop at the center.

Every non-loop cycle passes through the center.  A cycle cover
therefore chooses one path from the empty set to the full set
in the original subset lattice, and loops at all unused vertices.
Such a path adds the elements in the order of a unique permutation
\(\sigma\in S_m\), and its variable product is
\(\prod_{i=1}^m x_{i,\sigma(i)}\).
The merged path is a cycle of length \(m\), whose permutation
sign is \((-1)^{m-1}\).  The sign placed on its first edge
cancels that sign.  The \(s-m\) unused loops supply
\(L^{s-m}\).  The permutation expansion of the adjacency-matrix
determinant thus gives \eqref{eq:F009}, including its signs.
This is an integral polynomial identity.

Appending the block \(L I_{n-s}\) gives a matrix of size \(n\)
with determinant \(L^{n-m}\operatorname{perm}_m(X)\).
Thus the padded permanent itself is an endomorphism substitution
into \(\det_n\), and lies in \(\Omega_n\).  The latter is closed
and \(\mathrm{GL}(V)\)-stable, so it contains the whole
padded-permanent orbit closure.  Restriction of coordinate rings
is a surjection, and characteristic-zero semisimplicity gives
the inequality of every isotypic multiplicity.
\end{proof}

For \(m=1\) a diagonal matrix gives the assertion directly.
The identity remains valid after every substitution for \(L\),
including \(L=0\); when \(m\geq2\), \(s>m\), so both sides
then vanish. In particular one may put \(L=x_{11}\), the internal
padding convention of the supplied IP and BIP papers and of
\(q^m_\lambda(d[n])\) here. The construction is a morphism over the integers
from coefficient data to linear matrix pencils.  Its characteristic-zero
image containment and the induced coherent-algebra quotient are
the scheme and \(K_0\) meanings of the multiplicity statement.
The finite subset lattice describes the vertices, but its
permanent-specific edge labels are additional data; the containment
is not inferred from a formal lambda identity alone.

\paragraph{Exact small-size consequences in the two stated GCT ranges.}
For integer \(m\geq2\),
\[
 2^m-1\leq3m^4+1\quad(2\leq m\leq18),\qquad
 2^m-1\leq m^{25}\quad(2\leq m\leq189).
\]
Therefore an actual multiplicity obstruction in the range
\(n>3m^4\) must have \(m\geq19\); one in the range
\(n\geq m^{25}\) must have \(m\geq190\).
These bounds are combined with, and independent of, the row and
degree exclusions in M-008.  They are uses of a known determinantal
upper bound, not lower bounds on the complexity of the permanent.

Here is a direct verification that the cutoffs are exact for this
particular comparison.  The function
\(x\log2-c\log x\) is convex for \(x>0\).
At the endpoints \(m=2,18\), one has \(2^m\leq3m^4\);
convexity gives the inequality throughout that interval.
Likewise \(2^2\leq2^{25}\) and the exact integer comparison
\(2^{189}\leq189^{25}\) imply the second inequality
throughout \([2,189]\).
At the next integers, exact arithmetic gives
\[
 2^{19}-1=524287>390964=3\cdot19^4+1,
 \qquad 2^{190}-1>190^{25}.
\]
The ratios \((1+1/m)^4\) and \((1+1/m)^{25}\) decrease with
\(m\), and exact integer checks give
\(20^4<2\cdot19^4\) and \(191^{25}<2\cdot190^{25}\).
It follows inductively that the respective failures persist for
all larger integers, even with the displayed additive constants.
The complete integer comparisons and small determinant checks
are recorded in \texttt{permanent\_gate\_certificate.json}.

\subsection*{Finding F-010: the exact multiplicity-space datum left to compare}
Let \(V\) be a complex vector space, \(f\in\operatorname{Sym}^nV^*\),
and \(X_f=\overline{\mathrm{GL}(V)\cdot f}\).
For degree \(d\), use the natural decomposition
\[
 F_d(V)=\operatorname{Sym}^d\operatorname{Sym}^nV
       =\bigoplus_{\lambda\vdash nd}S_\lambda V\otimes M_\lambda,
 \qquad M_\lambda=[\lambda]^{S_n\wr S_d}.
\]
The \(\lambda\)-component of the evaluation tensor \(f^{\otimes d}\)
is an element of
\(S_\lambda V^*\otimes M_\lambda^*\), and therefore a linear map
\[
                    T_{f,\lambda}:M_\lambda\longrightarrow S_\lambda V^*.
\]
Then
\begin{equation}\label{eq:F010}
 \operatorname{mult}_{S_\lambda V}\mathbb C[X_f]_d
                  =\operatorname{rank}T_{f,\lambda}.
\end{equation}
In particular the remaining actual GCT comparison is exactly
\[
 \operatorname{rank}T_{L^{n-m}\operatorname{perm}_m,\lambda}
                 \ \mathop{\leq}^{?}\
 \operatorname{rank}T_{\det_n,\lambda}.
\]
It is a comparison of positive ranks, not merely the question whether
either map is zero. Here \(L=x_{11}\) for the stated GCT convention;
the rank characterization itself applies to every form and every
choice of padding.

\begin{proof}
Write
\(I(X_f)_d=\bigoplus_\lambda S_\lambda V\otimes J_{f,\lambda}\).
For \(a\in M_\lambda\), the condition
\(T_{f,\lambda}(a)=0\) says exactly that every polynomial
\(s\otimes a\), \(s\in S_\lambda V\), vanishes at \(f\).
Since this entire copy of \(S_\lambda V\) is stable under the
general linear group, it then vanishes on the entire orbit:
evaluation at \(g f\) is evaluation of \(g^{-1}(s\otimes a)\)
at \(f\).  It vanishes on the closure as well.
Conversely, a copy contained in the ideal vanishes at \(f\),
so its multiplicity vector is in \(\ker T_{f,\lambda}\).
Thus \(J_{f,\lambda}=\ker T_{f,\lambda}\).  Subtraction from
\(\dim M_\lambda\) proves \eqref{eq:F010}.
\end{proof}

The Young-set and wreath-invariant constructions determine the
ambient space \(M_\lambda\) and its dimension.  They do not
determine the map \(T_{f,\lambda}\), which depends on the
specific tensor \(f^{\otimes d}\).  One sufficient additional
statement would be a factorization
\[
 T_{L^{n-m}\operatorname{perm}_m,\lambda}
                   =A_\lambda T_{\det_n,\lambda}.
\]
This is equivalent to inclusion of the corresponding kernels
and is stronger than a rank inequality alone.  M-006 supplies
precisely such kernel inclusion in its stated projection range;
no such inclusion is established in general.

\paragraph{A small model of the missing information.}
The distinction is already visible for the standard representation
\(S=\mathbb C^2\) of \(\mathrm{GL}_2\), with
\(W=S\otimes\mathbb C^2\).  Regard elements of \(W\) as
two-column matrices.  The orbit closure of
\(v_1=(e_1,0)\) is the linear subspace \(S\otimes\mathbb C e_1\),
whereas the orbit of \(v_2=(e_1,e_2)\) is the invertible matrices
and its closure is all of \(W\).  In their degree-one coordinate
rings, the multiplicity of \(S^*\) is respectively one and two.
It is positive in both.  The ambient representation and all its
lambda operations are identical in the two comparisons.
The asserted orbit closures follow directly from transitivity
on nonzero vectors and density of invertible matrices, so this
example needs no external orbit classification.

This example refutes an inference from the ambient lambda object
and occurrence alone to equality, or to a chosen inequality, of
orbit-closure multiplicities.  It is not a counterexample to any
specific permanent-versus-determinant inequality.
Geometrically \(T_{f,\lambda}\) is a fiberwise linear map on a
Schur isotypic multiplicity bundle, whose minors cut out its
rank loci.  Projectors and the Schur decomposition are used in
characteristic zero, as in the polynomial-functor dictionary of
Macdonald, I, Appendix~A, (5.3)--(5.6), pp.~155--156.
The specific evaluation tensor, or equivalent equations for
these rank loci, is the additional datum not supplied by the
class in \(K_0\).

\subsection*{Finding F-011: a complete scalar-multiplicity family}
Let \(d\geq2\), \(n\geq1\), and \(0\leq k\leq n\). Delete zero
parts when forming the partition
\[
 \lambda_{d,k}=
       \bigl(d(n-k)+k,\underbrace{k,\ldots,k}_{d-1}\bigr)
                  \vdash nd.
\]
Then the ambient plethysm coefficient is
\begin{equation}\label{eq:F011}
 a_{\lambda_{d,k}}(d[n])=
 \begin{cases}1&k\text{ even},\\0&k\text{ odd}.
 \end{cases}
\end{equation}
As usual this is the stable coefficient of the Schur function;
when applying it to a fixed vector space, the partition must fit
its dimension. For \(d=1\), the partition is just \((n)\) and the
coefficient is one, independently of the redundant parameter \(k\).
The proof is a direct Cauchy-formula calculation; no historical
priority is claimed for the formula.

\begin{proof}
Put \(E=\mathbb C^d\) and take \(V\) of dimension at least \(d\).
In \(\operatorname{Sym}(V\otimes E)\), the subspace of column
multidegree \((n,\ldots,n)\) is naturally
\((\operatorname{Sym}^n V)^{\otimes d}\).
The permutation matrices in \(\mathrm{GL}(E)\) permute its
\(d\) tensor factors. Their invariant space is
\(\operatorname{Sym}^d\operatorname{Sym}^n V\).
The Cauchy decomposition
\[
 \operatorname{Sym}(V\otimes E)
      =\bigoplus_\mu S_\mu V\otimes S_\mu E
\]
therefore identifies the required multiplicity with
\[
 \dim\left((S_{\lambda_{d,k}}E)_{(n,\ldots,n)}\right)^{S_d}.
\]
The precise classical identity is Macdonald, Chapter~I,
equation~(4.3), printed p.~63, together with the polynomial-functor
interpretation in Appendix~A, (5.3)--(5.6), pp.~155--156.

Subtracting \(k\) from each of the \(d\) highest-weight coordinates
gives
\[
 S_{\lambda_{d,k}}E
       =\det(E)^k\otimes\operatorname{Sym}^{d(n-k)}E.
\]
Its weight \((n,\ldots,n)\) space is one-dimensional:
after the determinant twist it is spanned by
\(e_1^{n-k}\cdots e_d^{n-k}\).
Every permutation matrix fixes this monomial and acts on
\(\det(E)^k\) by \(\operatorname{sgn}^k\).
For \(d\geq2\), this one-dimensional representation has an
invariant exactly when \(k\) is even. This proves \eqref{eq:F011}.
The cases \(k=0\), \(k=n\), and \(d=1\) are included or handled
as stated, with no negative exponent.
\end{proof}

\paragraph{Consequences for actual and sufficient multiplicity tests.}
For \(d\geq2\), odd \(k\) gives zero permanent-side multiplicity
for every \(m,n\), since the ambient multiplicity is zero.
For even \(k\), there is no actual multiplicity obstruction in
the B\"urgisser--Ikenmeyer--Panova range \(n\geq m^{25}\).
Indeed \(q\leq a=1\), and their Theorem~1.4 says that
\(q>0\) implies \(K_n>0\); in that case \(q=1\leq K_n\).
For the Ikenmeyer--Panova range \(n>3m^4\), their Theorem~1.6
gives the analogous exclusion of the sufficient test \(q>g\),
since \(q>0\) forces \(g>0\). The latter argument alone does
not exclude an actual inequality \(q>K_n\).

More generally, throughout the BIP range every actual multiplicity
obstruction must have
\[
                     q^m_\lambda(d[n])\geq2,
                \qquad a_\lambda(d[n])\geq2.
\]
Thus all ambient multiplicity-one shapes are excluded there,
and \eqref{eq:F011} provides an explicit infinite family to which
this applies. No general determinant-side nonvanishing is asserted
for even \(k\) outside the stated no-occurrence range.
The special case \(k=2\) has a stronger all-\(n\) determinant
and permanent interpretation through the Hessian calculation in O-011.

\paragraph{Grothendieck dictionary.}
The multidegree condition is a character-weight submodule for the
diagonal torus in \(\mathrm{GL}(E)\); the permutation matrices
form its constant finite subgroup \(S_d\).
The calculation concerns the normalizer's action on a one-dimensional
weight submodule of a Schur functor. It uses locally free symmetric,
determinant, and Schur operations, then finite-group invariants in
characteristic zero. Thus it is a calculation in
\(K_0(\operatorname{Rep}(\mathrm{GL}(V)\times(T\rtimes S_d)))\),
with the grading/comodule dictionary of SGA~3, I, Proposition~4.7.3,
and the lambda operations of SGA~6, VI, Theorem~3.3.
The sign is the orientation character of the determinant line;
forgetting it would incorrectly predict the odd cases to occur.

\subsection*{Finding F-012: mixed near-rectangular tails have zero ambient multiplicity}
Let \(d\geq2\), \(n\geq1\), \(0\leq k<n\), and
\(0\leq s\leq d-1\). Put \(a=n-k\geq1\), \(D=da\), and
delete zero parts from
\[
 \lambda_{d,k,s}
  =\bigl(D+k-s,(k+1)^s,k^{d-s-1}\bigr)\vdash nd.
\]
This is a partition: \(D-s\geq1\), and for \(s>0\) its
first part is at least \(k+1\). Its ambient coefficient is
\begin{equation}\label{eq:F012}
 a_{\lambda_{d,k,s}}(d[n])=
 \begin{cases}
 1& s=0\text{ and }k\text{ even},\\
 1& s=d-1\text{ and }k\text{ odd},\\
 0&\text{otherwise}.
 \end{cases}
\end{equation}
The second positive endpoint is exactly F-011 with \(k\) replaced
by \(k+1\). In particular every mixed tail, with
\(1\leq s\leq d-2\), has zero permanent-side multiplicity,
for all sizes: its ambient multiplicity already vanishes.
At \(k=0\), the assertion recovers hook vanishing in the
indicated range.

\begin{proof}
Use the Cauchy weight-space description from F-011, with
\(E=\mathbb C^d\). Subtracting the rectangle \((k^d)\) gives
\[
 S_{\lambda_{d,k,s}}E
       =\det(E)^k\otimes S_{(D-s,1^s)}E.
\]
It remains to understand the balanced weight \((a,\ldots,a)\)
of the hook Schur module as a representation of the column
permutation group \(S_d\).

The hook form of Jacobi--Trudi gives, in the representation ring,
\[
 [S_{(D-s,1^s)}E]
    =\sum_{t=0}^{s}(-1)^{s-t}
       [\operatorname{Sym}^{D-t}E\otimes\Lambda^t E].
\]
This is the usual Jacobi--Trudi identity of Macdonald, I, (3.4),
expanded along the hook; no positivity is being inferred from
its signs. Taking the balanced weight commutes with restriction
to its normalizer, so it preserves this equality as an equality
of \(S_d\)-representation classes.

In the summand indexed by \(t\), a basis of that weight space is
indexed by the \(t\)-subsets \(J\subseteq[d]\): its basis vector is
\[
  \left(\prod_{i=1}^d e_i^{a-\mathbf1_{i\in J}}\right)
                 \otimes\bigwedge_{j\in J}e_j,
\]
with increasing wedge order. Every exponent is nonnegative
because \(a\geq1\). A column permutation permutes the first
factor and has exactly the oriented-subset sign in the second.
Consequently this weight space is the \(S_d\)-module
\(\Lambda^t P_d\), where \(P_d\) is the permutation module.
Write \(P_d=\mathbf1\oplus U_d\), with \(U_d\) the augmentation
representation. The exterior-series identity gives
\[
 \sum_{t=0}^s(-1)^{s-t}[\Lambda^t P_d]
                              =[\Lambda^s U_d].
\]
Thus the balanced-weight hook module has exactly that character.
Characteristic-zero semisimplicity identifies its isomorphism
class with \(\Lambda^s U_d\).

Including the determinant twist, the desired coefficient is
\[
                    \dim(\operatorname{sgn}^k\otimes
                                      \Lambda^s U_d)^{S_d}.
\]
The only invariant in \(\Lambda^\bullet U_d\) is in degree zero,
as proved by the orientation calculation in O-005.
Moreover \(U_d\) is self-dual and
\(\det(U_d)=\operatorname{sgn}\), so exterior duality gives
\[
  \operatorname{sgn}\otimes\Lambda^s U_d
                      \simeq\Lambda^{d-1-s}U_d.
\]
For even \(k\) the invariant occurs exactly at \(s=0\);
for odd \(k\), exactly at \(s=d-1\). Its dimension is one
in each case. This proves \eqref{eq:F012} and all stated
vanishing consequences.
\end{proof}

\paragraph{Dictionary and implication.}
This proof takes the torus-weight submodule of a Schur operation
on a locally free rank-\(d\) sheaf, then the invariant module
of the constant group scheme \(S_d\). Its exact computation
uses the exterior lambda identity for
\(P_d=\mathbf1\oplus U_d\), with the orientation character
retained. These are the grading and representation dictionaries
of SGA~3, I, 4.7.2--4.7.3, and the exact-sequence exterior
filtration of SGA~6, V, Lemma~2.2.1 and VI, Theorem~3.3.
It provides a genuine all-size exclusion of specified GCT
candidate shapes. It compares no positive multiplicities for
the shapes outside this vanishing family.

\subsection*{Finding F-013: exact ambient multiplicities after a rectangular shift}
Let \(d\geq2\), \(n\geq1\), \(0\leq k\leq n\), put
\(a=n-k\), \(D=da\), and let \(0\leq t\leq D/2\) be an integer.
Delete zero parts from
\[
                 \lambda=(D-t+k,t+k,k^{d-2})\vdash nd.
\]
Let \(p_j(d,b)\) count partitions of \(j\) in a rectangle with
\(d\) rows and width \(b\); set it equal to zero for \(j<0\)
or \(b<0\), and set \(p_0(d,0)=1\).
With \(\delta=d(d-1)/2\), one has
\begin{equation}\label{eq:F013}
 a_\lambda(d[n])=
 \begin{cases}
 p_t(d,a)-p_{t-1}(d,a),&k\text{ even},\\
 p_{t-\delta}(d,a-d+1)-p_{t-1-\delta}(d,a-d+1),
       &k\text{ odd and }a\geq d-1,\\
 0,&k\text{ odd and }a<d-1.
 \end{cases}
\end{equation}
Thus this ambient multiplicity is reduced to the same partition
differences as the rectangular/two-row Kronecker family F-002.
The formula is a direct consequence of the Cauchy and two-row
Jacobi--Trudi identities, with the sign action recorded explicitly.
It is not a claim that the two orbit closures have these full
ambient multiplicities.

\begin{proof}
Again put \(E=\mathbb C^d\) and use the column weight
\((n,\ldots,n)\) description in F-011. Removing the determinant
twist gives
\[
 S_\lambda E=\det(E)^k\otimes S_{(D-t,t)}E.
\]
Jacobi--Trudi, equivalently the difference of the two adjacent
horizontal Pieri decompositions, is
\[
 [S_{(D-t,t)}E]
 = [\operatorname{Sym}^{D-t}E\otimes\operatorname{Sym}^{t}E]
  -[\operatorname{Sym}^{D-t+1}E\otimes
                                      \operatorname{Sym}^{t-1}E],
\]
with the second summand zero for \(t=0\).
The classical locators are Macdonald, I, (3.4) and (5.16).

The balanced weight \((a,\ldots,a)\) in the first tensor product
has a monomial basis indexed by
\[
             B_t=\{\beta\in\{0,\ldots,a\}^d:|\beta|=t\}:
 \qquad e^{a\mathbf1-\beta}\otimes e^\beta.
\]
The column permutation group permutes this basis. The second
tensor product gives \(\mathbb C[B_{t-1}]\) in exactly the same
way. Thus the required coefficient is the difference of the
\(\operatorname{sgn}^k\)-multiplicities in those two permutation
modules.

If \(k\) is even, each orbit contributes one invariant.
Sorting the entries of \(\beta\) identifies its orbit with a
partition of \(t\) having at most \(d\) parts, each at most \(a\).
This gives \(p_t(d,a)-p_{t-1}(d,a)\).

If \(k\) is odd, an orbit contributes to the sign isotypic
space precisely when its stabilizer has no odd permutation.
The stabilizer is the product of the symmetric groups on the
equal-entry positions. It contains an odd transposition if any
two entries are equal, and is trivial otherwise. Hence the
contributing orbits are exactly the strictly increasing lists
\[
                0\leq b_1<b_2<\cdots<b_d\leq a.
\]
There are none when \(a<d-1\). Otherwise subtract
\((0,1,\ldots,d-1)\). The resulting list is weakly increasing,
lies in \([0,a-d+1]\), and has sum \(t-\delta\).
Reversing its order is a partition in the stated rectangle.
Applying the same bijection at \(t-1\) proves the second line
of \eqref{eq:F013}.
All finite-group invariant and sign-isotypic functors are exact
in characteristic zero, so the subtraction is legitimate.
\end{proof}

In the odd case, the nonzero shifted degree \(u=t-\delta\)
satisfies \(u\leq d(a-d+1)/2\). Thus the displayed differences
stay in the first half of their rectangles, consistently with
their representation-theoretic nonnegativity. The conventions
also cover \(a=0,t=0\), and the single-point rectangle in the
odd case \(a=d-1,t=\delta\).

\paragraph{A concrete four-copy family beyond occurrence.}
If \(d\geq6\) and \(n\geq8\), take \(k=2,t=6\). Then
\[
 \lambda=(d(n-2)-4,8,2^{d-2}),\qquad
                         a_\lambda(d[n])=4.
\]
Indeed the rectangle bounds do not affect partitions of six or
five. Removing the part of size one leaves the four partitions
\((6),(4,2),(3,3),(2,2,2)\).
For \(d\leq m^2\), \(m\geq3\), and \(d\geq6\), this family
also meets the coarse internal-padding conditions:
its tail has size \(2d+4\leq md\) and its length is \(d\).
The ambient coefficient four does not prove positive permanent
multiplicity or compare it with the determinant closure.
It gives a precise small multiplicity space on which F-010's
remaining rank comparison can be studied, including at arbitrarily
large row counts.

\paragraph{Dictionary.}
This is a balanced-torus-weight calculation of a Schur functor
followed by the trivial or orientation isotypic summand for
the constant group scheme \(S_d\). The two permutation modules
are linearizations of the actual finite sets \(B_t\) and
\(B_{t-1}\). The sign contribution is detected by the stabilizer,
so repeated entries are excluded for a representation-theoretic
reason, not by an assumption that a virtual Burnside class is
effective. The character-weight and comodule language is SGA~3,
I, 4.7.2--4.7.3; symmetric, determinant and Schur operations
are those of SGA~6, VI, Theorem~3.3 and the polynomial-functor
dictionary already used above.
